% Bibliography database embedded directly in this .tex file.
% The filecontents environment writes apssamp.bib at compile time, so no
% external .bib file is needed. Edit the entries below to manage references.
% (Example entries -- replace with your own as needed.)
\begin{filecontents}[overwrite]{apssamp.bib}
@article{example2024,
  author  = {J. Doe and A. Smith},
  title   = {An example reference for the bibliography},
  journal = {Phys. Rev. Lett.},
  volume  = {132},
  pages   = {010101},
  year    = {2024},
  doi     = {10.1103/PhysRevLett.132.010101},
}
\end{filecontents}
% ****** Start of file apssamp.tex ******
%
%   This file is part of the APS files in the REVTeX 4.2 distribution.
%   Version 4.2a of REVTeX, December 2014
%
%   Copyright (c) 2014 The American Physical Society.
%
%   See the REVTeX 4 README file for restrictions and more information.
%
% TeX'ing this file requires that you have AMS-LaTeX 2.0 installed
% as well as the rest of the prerequisites for REVTeX 4.2
%
% See the REVTeX 4 README file
% It also requires running BibTeX. The commands are as follows:
%
%  1)  latex apssamp.tex
%  2)  bibtex apssamp
%  3)  latex apssamp.tex
%  4)  latex apssamp.tex
%
\documentclass[%
 reprint,
%superscriptaddress,
%groupedaddress,
%unsortedaddress,
%runinaddress,
%frontmatterverbose, 
%preprint,
%preprintnumbers,
%nofootinbib,
%nobibnotes,
%bibnotes,
 amsmath,amssymb,
 aps,
%pra,
%prb,
%rmp,
%prstab,
%prstper,
%floatfix,
]{revtex4-2}

\usepackage{graphicx}% Include figure files
\usepackage{dcolumn}% Align table columns on decimal point
\usepackage{bm}% bold math
\usepackage{float}
%\usepackage{hyperref}% add hypertext capabilities
%\usepackage[mathlines]{lineno}% Enable numbering of text and display math
%\linenumbers\relax % Commence numbering lines

%\usepackage[showframe,%Uncomment any one of the following lines to test 
%%scale=0.7, marginratio={1:1, 2:3}, ignoreall,% default settings
%%text={7in,10in},centering,
%%margin=1.5in,
%%total={6.5in,8.75in}, top=1.2in, left=0.9in, includefoot,
%%height=10in,a5paper,hmargin={3cm,0.8in},
%]{geometry}

\begin{document}

\preprint{APS/123-QED}

\title{LensAFM}

\author{Minwoo Kim (Team 1)}

\date{\today}% It is always \today, today,
             %  but any date may be explicitly specified

\begin{abstract}
An article usually includes an abstract, a concise summary of the work
covered at length in the main body of the article. 
\end{abstract}

\maketitle

%\tableofcontents

\section{\label{sec:level1}Section 1}

This is section \ref{sec:level1}. Cite like this \cite{example2024}.

\subsection{Subsection}

\subsubsection{Subsubsection}

\section{Figures and Tables}

Use the \texttt{figure} environment for a figure, the \texttt{table} environment for a table.
\begin{figure}[H]
\centering
\includegraphics[width=.7\linewidth]{images/C_short_leveled}
\caption{Caption}
\end{figure}

In PRL/REV\TeX\ style the caption sits \emph{above} the table, and the
\texttt{ruledtabular} environment draws the top/bottom Scotch rules with a
\texttt{\textbackslash colrule} under the header. See Table~\ref{tab:sample}.
\begin{table}[H]
\caption{\label{tab:sample}%
Measured rod height and FWHM for the Contact (C) and Non-contact (NC) tips.}
\begin{ruledtabular}
\begin{tabular}{lcc}
Tip & Height (nm) & FWHM (nm) \\
\colrule
Contact     & $67.6 \pm 1.2$ & $148 \pm 6$ \\
Non-contact & $52.3 \pm 0.9$ & $131 \pm 5$ \\
\end{tabular}
\end{ruledtabular}
\end{table}

\begin{acknowledgments}
Talk about LLM assistance.
\end{acknowledgments}

\appendix

\section{Author Contributions}
\begin{itemize}
    \item Minwoo Kim: Wrote this report. Also, ...
    \item Heemin Kim: ...
    \item Seokhun Oh: ...
    \item Chihoon Choi: ...
\end{itemize}


\bibliography{apssamp}% Produces the bibliography via BibTeX.

\end{document}
%
% ****** End of file apssamp.tex ******
